\documentclass[10pt,a4paper]{amsart}
\usepackage[margin=19mm]{geometry}
\usepackage{amsmath,amssymb,mathtools}
\usepackage[scr=rsfs,cal=euler]{mathalfa}
\usepackage{xfrac}
\usepackage[unicode,hidelinks,bookmarks=false]{hyperref}
\newcommand{\ie}{i.e.\ }
\newcommand{\cf}{cf.\ }
\newcommand{\resp}{respectively\ }
\newcommand{\Subsection}{\subsection}
\providecommand{\nobreakdash}{\nobreak\hbox{-}}
\setlength{\emergencystretch}{2em}
\multlinegap0pt
\usepackage{amssymb}
\usepackage{enumerate}
\usepackage[normalem]{ulem}
\usepackage{cancel}
\usepackage{mathtools}
\usepackage{bm}
\usepackage{xcolor}
\newtheorem{lemma}{Lemma}
\newcommand{\R}{\mathbb{R}}
\newcommand{\cX}{\mathcal{X}}
\title{Overdamped limits for Langevin dynamics with position-dependent coefficients via $L^2$-hypocoercivity}
\author{No\'e Blassel}
\date{Source: arXiv:2602.16924v3 (CC BY 4.0)}
\begin{document}
\maketitle

\noindent\emph{Source and licence.} Overdamped limits for Langevin dynamics with position-dependent coefficients via $L^2$-hypocoercivity. Version arXiv:2602.16924v3 (CC BY 4.0), distributed by arXiv under the Creative Commons Attribution 4.0 International licence, http://creativecommons.org/licenses/by/4.0/, as recorded in the arXiv licence metadata for this exact version and shown on the abstract page https://arxiv.org/abs/2602.16924v3. Full licence notice: \texttt{licenses/arxiv-2602.16924-CC-BY-4.0.txt}.\par\smallskip

\noindent\textit{Editorial scope.} Counted unit: the article's lemma lemma:commutation (source0.tex lines 1441--1444). The article's complete original proof of it is delivered with it (source0.tex lines 1445--1457, one \texttt{proof} environment of 13 lines). No mathematics is written, repaired or re-derived: the statement and the proof are the article's own lines, in the article's own notation, and no retained line is substituted or re-flowed.  Retained uncounted as the source-local context the proof needs: lines 975--978.  Nothing was omitted from any retained span, so the package carries no omission record.  No retained line contains a citation, so the wrapper carries no bibliography and no bibliography entry.  No retained line contains a figure environment, an image-inclusion command or drawing code, so no image is delivered and the package declares no asset dependency.  Editorial formatting only, in addition: an AMS \texttt{amsart} preamble that restates the article's own declarations and macros used by the retained lines, namely the environments \texttt{lemma} on separate counters, together with the standard packages those lines need.  The retained environments print in this document as Lemma 1 in order of appearance, whereas the article prints the same items with the numbering of its own full document.  The complete native source of the article as distributed in the arXiv e-print for this exact version is bundled unchanged as \texttt{sources/194931/source0.tex}.
\begin{equation}
    \label{eq:symmetry_condition}
\nabla_q x^\top \nabla_q v = A^{-1}(q)\nabla_q v = \nabla_q v^\top \nabla_q x = \nabla_q v^\top A^{-\top},
\end{equation}

\begin{lemma}[Index computation]
    \label{lemma:commutation}
    Let~$A:\cX\to\R^{d\times d}$ be a smooth matrix field such that~$A^{-\top}$ has gradient rows, and~$v(q,p)=A(q)p$. Then the relation~\eqref{eq:symmetry_condition} holds.
\end{lemma}

\begin{proof}
    Since~$A^{-1}$ has gradient columns, we may write locally~$(A^{-1})_{ij} = \partial_i \phi_j$ for some~${(\phi_j)_{1\leq j\leq d}\in\mathcal C^\infty(\cX)^d}$.
    We compute, with derivatives taken with respect to the~$q$-variable:
    \begin{equation}
        \begin{aligned}
            (A^{-1}\nabla_q v)_{ij} &= \sum_{\alpha,\beta} (A^{-1})_{i\alpha}\partial_j(A_{\alpha\beta})p_\beta\\
            &=-\sum_{\alpha,\beta}(A^{-1})_{i\alpha}\left[A\partial_j(A^{-1})A\right]_{\alpha\beta}p_\beta\\
            &=-\sum_{\gamma,\beta}\partial_j(A^{-1})_{i\gamma}A_{\gamma\beta}p_\beta\\
            &=-\sum_{\gamma}\partial^2_{ji}\phi_\gamma v_\gamma,
        \end{aligned}
    \end{equation}
    using the matrix identity~$\partial_j A = -A\partial_j(A^{-1})A$. This shows that~$A^{-1}\nabla_q v$ is symmetric, and therefore that~\eqref{eq:symmetry_condition} holds.
\end{proof}

\end{document}
